\section{Give the sprite a file representation}
The eight-byte sprite stored one on-or-off decision per bit. Its eight row bytes were
\texttt{16 3c 7e db ff 24 42 81}, in hexadecimal, with bit seven at the left edge of
each row. Keep that image. Expand each set bit into an opaque white pixel and each
clear bit into an opaque black pixel. There are now sixty-four pixels, each with four
components: red, green, blue and alpha. Alpha records opacity; 255 is fully opaque and
zero is fully transparent.

QOI, the Quite OK Image format, describes an image as a succession of commands that
reconstruct pixels. A command can supply a complete color, change the previous color,
retrieve a cached color, or repeat it. The decoder has to retain state between commands.
This differs from PPM, where each RGB triplet describes its pixel without referring to
an earlier triplet.

The project has three useful boundaries. First recover exact RGBA values in memory.
Then connect that memory to the display and observe the order in which pixels are
written. Finally, encode those values into a file that another decoder can open. A
recognizable sprite tests orientation, but exact component comparisons are still
needed: a wrong alpha value or a cache collision can escape a visual check.

The fixture helper in \path{drafts/qoi-fixtures} writes the sprite or a smaller worked
trace. Run \texttt{python3 fixture.py new-image.qoi} for the sprite. The helper refuses
an existing output path. Its alternate PPM output is an RGB preview and discards alpha;
it is not an alpha-preserving conversion.
